\documentclass[11pt]{article}
\usepackage{preamble}
\setlength{\parskip}{4pt}

\begin{document}

\daybanner{AP Statistics \textperiodcentered\ Unit 4 \textperiodcentered\ Topic 4.9 \textperiodcentered\ B: 2027-01-28 \textperiodcentered\ E: 2027-03-17 \textperiodcentered\ TEACHER KEY}{Two Test Plans for Cure Time}

\sectionbanner{CED TETHER}

{\small
\begin{itemize}[leftmargin=1.5em,itemsep=2pt,topsep=2pt]
  \item \textbf{Skill 1.E} --- Identify an appropriate inference method for significance tests.
  \item \textbf{Skill 1.F} --- Identify null and alternative hypotheses.
  \item \textbf{Skill 4.C} --- Verify that inference procedures apply in a given situation.
  \item \textbf{EU VAR-7} --- The t-distribution may be used to model variation.
  \item \textbf{LO VAR-7.F} --- Identify an appropriate selection of a testing method for a difference of two population means. [Skill 1.E]
  \item \textbf{EK VAR-7.F.1} --- For a quantitative variable, the appropriate test for a difference of two population means is a two-sample t-test for a difference of two population means.
  \item \textbf{LO VAR-7.G} --- Identify the null and alternative hypotheses for a difference of two population means. [Skill 1.F]
  \item \textbf{EK VAR-7.G.1} --- The null hypothesis for a two-sample t-test for a difference of two population means, $\mu_1$ and $\mu_2$, is: $H_0$: $\mu_1$ - $\mu_2$ = 0, or $H_0$: $\mu_1$ = $\mu_2$. The alternative hypothesis is $H_a$: $\mu_1$ - $\mu_2$ < 0, or $H_a$: $\mu_1$ - $\mu_2$ > 0, or $H_a$: $\mu_1$ - $\mu_2$ $\neq$ 0, or $H_a$: $\mu_1$ > $\mu_2$, or $H_a$: $\mu_1$ < $\mu_2$, or $H_a$: $\mu_1$ $\neq$ $\mu_2$.
  \item \textbf{LO VAR-7.H} --- Verify the conditions for the significance test for the difference of two population means. [Skill 4.C]
  \item \textbf{EK VAR-7.H.1} --- In order to make statistical inferences when testing a difference between population means, we must check for independence and that the sampling distribution is approximately normal:
  \item a. Individual observations should be independent: (i) Data should be collected using simple random samples or a randomized experiment; (ii) When sampling without replacement, check that $n_1$ $\leq$ 10\% $N_1$ and $n_2$ $\leq$ 10\% $N_2$.
  \item \textbf{b. The sampling distribution of $\bar{x}_1$ - $\bar{x}_2$ should be approximately normal (shape): (i) If the observed distribution is skewed, both $n_1$ and $n_2$ should be greater than 30; (ii) If the sample size is less than 30, the distribution of the sample data should be free from strong skewness and outliers} --- checked for BOTH samples.
\end{itemize}
}
\textbf{Essential question:} How do the original question and study design control a defensible two-mean test plan?\par
\textbf{DOK-3 skill:} 4.C \textperiodcentered\ \textbf{FRQ pattern:} {\small\texttt{predata-\allowbreak{}tail-\allowbreak{}parameter-\allowbreak{}and-\allowbreak{}scope}}\par
\textbf{DOK rationale:} Part (c) coordinates the pre-data question, parameter order, sample structure, distribution evidence, and assignment limits to adjudicate two test plans and trace three misleading conclusions.\par

\frameworkphaseheader{Do Now (first take) $\rightarrow$ Explore (video + follow-along) $\rightarrow$ Exit (finish + turn in)}{1 $\rightarrow$ 3}{45}{%
\begin{itemize}[leftmargin=1.5em,itemsep=2pt,topsep=2pt]
  \item Display the board and ask for one sentence using the study description.
  \item Begin the 7.8 video follow-along at minute 5.
  \item Return to the board at minute 33. Students finish the shortened parts and justify their judgment.
  \item Collect at the door. Score part (c) only using the three required elements.
\end{itemize}
}{%
\begin{itemize}[leftmargin=1.5em,itemsep=2pt,topsep=2pt]
  \item Write a one-sentence first take before the video.
  \item Complete the video follow-along worksheet.
  \item Finish (a)--(c), revise the first take in the exit line, and turn in the sheet.
\end{itemize}
}{%
\begin{itemize}[leftmargin=1.5em,itemsep=2pt,topsep=2pt]
  \item Which specific number or study detail supports your choice?
  \item What claim would go beyond the evidence, and why?
\end{itemize}
}{Circulate and ask for evidence. Accept a concise response; do not require omitted calculations or a second full inference write-up.}

\begin{callout}{\IconWarn\ WATCH FOR}{calloutred}
\begin{itemize}[leftmargin=1.5em,itemsep=2pt,topsep=2pt]
  \item Choosing a speaker without a reason.
  \item Copying numbers without explaining what they support.
\end{itemize}
\end{callout}

\sectionbanner{SCORING PART (c) --- E / P / I}

\textbf{Expected elements}
\begin{itemize}[leftmargin=1.5em,itemsep=2pt,topsep=2pt]
  \item \textbf{predata-tail} (required) --- Uses the nondirectional original question to justify a two-sided alternative
  \item \textbf{parameters} (required) --- Explains why unknown population means belong in the hypotheses
  \item \textbf{cause} (required) --- Explains that formula and line differences cannot be separated without random assignment
\end{itemize}

{\small\renewcommand{\arraystretch}{1.25}
\noindent\begin{tabular}{|p{0.5in}|p{5.9in}|}
\hline
\textbf{E} & Justifies the original two-sided question, population notation, and causal limit. \\ \hline
\textbf{P} & Shows substantive valid reasoning for at least one required element, but does not meet all E requirements. Credit correct reasoning even when another claim is wrong. \\ \hline
\textbf{I} & Shows no substantive valid reasoning for any required element; a name, unsupported choice, or copied number alone does not count. \\ \hline
\end{tabular}}

\clearpage
\focusbanner{TODAY'S PROBLEM --- annotated}

Suppose line N uses a new formula for sealant, a material that closes gaps. Line S uses the standard formula. A team takes independent simple random samples of 18 of 600 recent N batches and 18 of 540 recent S batches. Mean cure times are 41.2 and 44.0 minutes; the boxplots are shown. Before sampling, the team asked, ``Do the two populations differ in mean cure time?'' Formula was not randomly assigned to lines.\par\smallskip \textbf{Rhea:} ``Keep the original question and use population means. A difference between lines may have other causes.''\par \textbf{Sol:} ``Since $41.2<44.0$, use a lower tail with sample means. A significant result proves the formula caused faster curing.''\par
\begin{center}
\begin{tikzpicture}
\begin{axis}[
  width=0.49\linewidth,
  height=1.15in,
  ytick={2,1}, yticklabels={{Line N},{Line S}}, ymin=0.4, ymax=2.6,
  axis x line=bottom, axis y line*=left, y axis line style={draw=none}, boxplot/box extend=0.5, xmajorgrids, enlarge x limits=0.06,
  xlabel={Cure time (min)}, tick label style={font=\small}, label style={font=\small},
  ytick style={draw=none}, yticklabel style={font=\small\bfseries}
]
\addplot[boxplot prepared={lower whisker=34, lower quartile=38.5, median=41, upper quartile=44, upper whisker=48.5, draw position=2}, fill=calloutblue, draw=dayblue] coordinates {};
\addplot[boxplot prepared={lower whisker=35.5, lower quartile=40, median=44, upper quartile=48, upper whisker=52.5, draw position=1}, fill=calloutblue, draw=dayblue] coordinates {};
\end{axis}
\end{tikzpicture}
\end{center}
\begin{firsttakebox}{FIRST TAKE (5 min)}
Did the team originally ask whether N was faster or whether the lines differed in either direction? Explain in one sentence.\par
\answer{Ungraded. Everyday language is enough before the video; formal vocabulary belongs in the finish.}
\end{firsttakebox}

\textit{Rules callout ``HYPOTHESES FOR TWO MEANS (keep on the page)'' is printed on the student sheet.}\par\medskip

\dokbadge{1}~\textbf{(a)} Define $\mu_N$ and $\mu_S$ in context.\par
\answer{$\mu_N$ and $\mu_S$ are the mean cure times, in minutes, for the 600 recent N batches and 540 recent S batches.}

\dokbadge{2}~\textbf{(b)} State $H_0$ and $H_a$ in N-minus-S order to match the question asked before sampling.\par
\answer{$H_0:\mu_N-\mu_S=0$; $H_a:\mu_N-\mu_S\neq0$.}

\dokbadge{3}~\textbf{(c)} In 3--4 sentences, judge the two plans. Explain why the original question sets the tail, why hypotheses use population means, and why these two lines cannot isolate the formula as the cause.\par
\answer{Rhea is justified because the original question asked about a difference in either direction, so the alternative is two-sided. Choosing a lower tail only after seeing $41.2<44.0$ changes that question to favor the observed result. Hypotheses concern unknown population means; the sample means are already known. Because formula was not randomly assigned to lines, other line differences could explain the gap, so a significant result would not establish formula causation.}

\begin{summaryexitbox}{EXIT}
One thing the video changed about my first take: \hrulefill
\end{summaryexitbox}

\end{document}
